\section{The completion problem}\label{sec:intro}
A terminal statistic may be inexpensive to store and impossible to maintain with the same online information. The missing object is not a dimension of the formal state space, but a compatibility condition between acquisition and every intermediate act of forgetting. This paper studies that condition for a prepared causal experiment and a specified decision task.

Our starting data are positive normalized kernels, legal interventions, and physically executable future tests. Their expectations identify histories that have the same predictive content. Mixtures of preparations give a convex predictive carrier, whereas an implemented acquisition policy gives a probability law on that carrier. These two objects are different. Compression replaces an acquired state by a conditional barycenter. It therefore contracts the acquired law in convex order, but a sequence of such contractions does not by itself describe an autonomous finite machine: the same persistent label cannot invoke different transition programmes at different unobserved times.

Theorem~\ref{thm:completion} gives an exact answer to both issues. For an externally clocked observer, a finite-support law and an action distribution at each retained state suffice; one tests a one-step barycentric inequality at each cut. For an autonomous observer, one keeps the proposed label occupancies and their conditional predictive states at every cut. A simultaneous dual inequality tests whether a single label/action/report transition rule can realize all of them. It eliminates the unknown transition kernels from the feasibility condition, without presuming a phase index or an exact posterior register. The criterion is variational, not a closed-form exponent or an efficient synthesis algorithm. A local converse certifies an infeasible proposed law sequence. Theorem~\ref{thm:finite_obstruction} additionally gives optimal external observers and finite test families excluding every proposed flow below an unattainable risk threshold. This global compactness conclusion has no automatic effective counterpart for general stationary Borel programmes.

The exact mechanism uses no acquired density, finite dimension, smooth task, summable stability, or common report support. Theorem~\ref{thm:completion} gives the risk frontier and Corollary~\ref{cor:zero} gives an exact zero-loss criterion. In contrast, a power-law conclusion needs quantitative hypotheses. Under an all-state/all-action density bound, block summability, and strong curvature of the task, Theorem~\ref{thm:regular} gives, uniformly in the horizon $n$,
\begin{equation}\label{eq:intro_regular}
 R_{\cp}(n,M)-B_n^*\asymp R_{\ext}(n,M)-B_n^*\asymp M^{-2/d},
 \qquad R_{\aut}(n,M)-B_n^*\asymp(M-n)^{-2/d}.
\end{equation}
The last formula assumes clock-neutral report supports and an internally exact deadline; it is not the conventional externally clocked filtering law. The all-state density condition is stronger than positive observation noise, and the strong-curvature condition is absent from the exact criterion.

A different sharp law makes the distinction substantive. When one capture reveals a prepared variable and all later reports are blind, Theorem~\ref{thm:blind} proves
\begin{equation}\label{eq:intro_blind}
 R_{\ext}(\mathbf m)=q_{\min_t m_t}(\mu),\qquad
 R_{\aut}(n,M)=q_{\lfloor(M-1)/n\rfloor}(\mu).
\end{equation}
Here $q_L(\mu)$ is the squared quantization error of the \emph{actually captured} distribution. The predictor must carry its resolution through every intervening cut; spending almost all labels at the terminal cut cannot restore discarded information. The acquired law may be fractal or countably atomic. This noncontracting class exhibits a different memory law without leaving the same general criterion.

The Gaussian and quantum realizations in Section~\ref{sec:realizations} establish the regular hypotheses from raw kernels for familiar probability forecasts: categorical filtering under Brier loss and prediction of a fixed informationally complete measurement. The latter is a finite-dimensional, known-instrument theorem; noncommutativity does not imply a result for arbitrary quantum control. The numerical statements give error propagation \emph{given} integration, selection and execution certificates, not polynomial-time synthesis or universal computability.

The one-step convex-order theorem is classical in substance\cite{strassen,blackwell}; conditional means, Bellman recursion and convex separation are not claimed as new. The contribution proposed here is their exact causal composition with counted finite supports, simultaneous phase-reuse constraints, an action-slack identity under the deployed law, and matching interface-dependent consequences. The mass--curvature and block-stable arguments are inherited and re-proved in the form used here\cite{prior}. Complete earlier articles are retained as archival companions, but no proof in this article requires a realization theorem or an unproved assertion from them.
